% !TEX program = pdflatex
\documentclass[11pt]{article}
\setlength{\textwidth}{6.5in}
\setlength{\oddsidemargin}{0pt}
\setlength{\evensidemargin}{0pt}
\setlength{\textheight}{9in}
\setlength{\topmargin}{0pt}
\addtolength{\topmargin}{-\headheight}
\addtolength{\topmargin}{-\headsep}
\ifdefined\pdfpagewidth \pdfpagewidth=8.5in \pdfpageheight=11in \fi
\ifdefined\pagewidth \pagewidth=8.5in \pageheight=11in \fi
\usepackage{amsmath,amssymb}
\usepackage{array}
\usepackage{booktabs}
\usepackage[tableposition=top]{caption}
\usepackage{longtable}
\usepackage{enumitem}
\usepackage{tikz}
\usepackage{pgfplots}
\pgfplotsset{compat=1.18}
\emergencystretch=3em

% ---- notation ---------------------------------------------------------------
\newcommand{\um}{\,$\mu$m}
\newcommand{\kms}{km\,s$^{-1}$}
\newcommand{\flux}{erg\,cm$^{-2}$\,s$^{-1}$}
\newcommand{\mA}{m\AA}
\newcommand{\sci}[2]{#1\times10^{#2}}
\newcommand{\ion}[2]{#1\,\textsc{#2}}
\newcommand{\nuc}[2]{$^{#1}$#2}
\newcommand{\lxl}{$L_X/L_\mathrm{bol}$}

% ---- front matter -----------------------------------------------------------
\title{Surfaces, Shapes, Spins and Orbits:\\ Archival Tests for Artifacts Among the Small Bodies}
\author{Charles C. Norton}
\date{}

\begin{document}
\maketitle

\begin{abstract}
An artifact among the asteroids, comets and meteoroids would most likely be found in the
archives that already hold its brightness, colour, shape, spin and orbit. We apply
thirteen tests to them. Of $315$ radar detections, every albedo above $0.3$ belongs to
a metal-rich or enstatite body and none is the specular return of a flat reflector.
Gaia DR3 reflectance spectra of $58208$ bodies hold no solar-cell surface, which
injections recover at $10$ per cent of a body's light in $88$ per cent of cases. In
$4723$ ALCDEF sessions no brightening recurs at a body's rotation period beyond chance.
No ATLAS phase curve is flat at a high albedo, and no NEOWISE fit stays hotter than
sunlight allows at every epoch. Reduced to zero phase, no lightcurve amplitude among
$31008$ records exceeds the $2.5$\,mag of a $10:1$ plate. The longest straight limbs in
$395$ stellar occultations lie on bodies whose shape models carry flat facets. No
kilometre body spins faster than ten minutes or needs more than $100$\,kPa to hold
together. The one triple of asteroids that share a period, magnitude and amplitude on
unrelated orbits is three S-types at the spin barrier's pile-up, a one-in-$25$ chance.
Among $301$ bodies with radial non-gravitational terms, three are lighter than a rubble
pile, each with a published explanation. No near-Earth object occupies the cells the
delivery model leaves empty in a way dynamics forbids. The densest hyperbolic meteor
group of the Global Meteor Network's $3.58$ million orbits is matched by a $1.5$ per
cent speed bias, and its focal line holds no stream. NELIOTA's $284$ lunar impact
flashes are thermal and recur in place no more often than chance.
\end{abstract}

% =============================================================================
\section{Introduction}\label{sec:intro}
% =============================================================================

The search for artifacts in the Solar System~\cite{Freitas1983,HaqqMisra2012} asks
whether something made, a probe, a reflector, a habitat or a fragment, is among the
bodies we already catalogue~\cite{Benford2019}. Surveys of the small bodies were not
designed for that question, but they hold what it needs: radar echoes, reflectance
spectra, lightcurves, phase curves, thermal fits, occultation chords, non-gravitational
terms in the orbit solutions, and the orbits of meteoroids and of the near-Earth
population. A made object can differ from a natural body in each: a flat, conducting or
photovoltaic surface; a shape no collisional fragment has; a spin no rubble pile
survives; an area too large for its mass; an orbit the natural sources do not reach.
This paper puts each of these to the archive that measures it, with a null drawn from
the same data, and states what each would have found.

\subsection{Contributions}
\begin{enumerate}[leftmargin=*,itemsep=1pt]
\item Tests of the surfaces of the small bodies: radar albedos, Gaia DR3 reflectance
spectra against metal and photovoltaic templates with injected surfaces, glints in
ALCDEF lightcurves, ATLAS phase curves and NEOWISE beaming parameters.
\item Tests of their shapes and spins: lightcurve amplitudes at zero phase, straight
limbs in occultation profiles, the strength that the fastest spins need, and bodies that
share a period, size and amplitude on unrelated orbits.
\item Tests of their orbits: area-to-mass ratios from non-gravitational terms, near-Earth
objects where the delivery model has none, and coherent hyperbolic meteor streams.
\item A test of NELIOTA's lunar impact flashes for non-thermal colours and recurrence.
\end{enumerate}

% =============================================================================
\section{Surfaces}\label{sec:surfaces}
% =============================================================================

\emph{Radar.} A flat metal plate facing the radar returns its echo in the opposite
circular sense with a radar albedo near one, where a rough natural surface returns some
power in the same sense and a stony one has an albedo near $0.1$~\cite{Shepard2015}. The
compilation of every published asteroid radar detection in the Small Bodies Node, $315$
objects, holds nine albedos above $0.3$: M- and X-class bodies, the radar-metallic
1986~DA~\cite{Ostro1991}, the E-types Nereus and 1998~WT24, and the Sq-type Bacchus at
$0.33\pm0.25$. None combines an albedo near one with a same-sense to opposite-sense ratio
near zero, the signature of a specular reflector.

\emph{Reflectance spectra.} Gaia DR3 gives reflectance spectra of $60518$ Solar System
objects in $16$ bands from $374$ to $1034$\,nm~\cite{Galluccio2023}, of which $58208$
have twelve or more good bands from $418$ to $990$\,nm. Each is fitted by the first
four principal components of the sample, with and without the reflectance of a
photovoltaic cell and of gold, copper, silver and aluminium from their optical
constants~\cite{Polyanskiy2024} as an additive component. No spectrum takes a
photovoltaic component, and the drops in $\chi^2$ above $25$ that the metal templates
give fall on V, A and S types whose deep absorption bands the four components fit
poorly. Injected into $3000$ spectra, a photovoltaic surface is recovered in $41$, $88$
and $100$ per cent of cases at $5$, $10$ and $20$ per cent of the light at $550$\,nm;
gold and copper in $19$ and $9$ per cent at $50$ per cent of the light, and silver and
aluminium, whose reflectance is nearly flat across the band, never.

\emph{Glints.} A flat surface on a spinning body reflects the Sun toward us once a
rotation, as a brief spike at a fixed rotational phase. In the ALCDEF photometry,
$4723$ sessions of $20$ or more points at $15$\,min or less on bodies with an LCDB
period~\cite{Warner2009} of quality $2$ or better, spikes above the rotational
lightcurve that recur at the body's period are compared with the same statistic for
fades. No body's spikes recur at $p<0.01$, where the fades, tested the same way, gather
in ten groups. Injected spikes are recovered in $3$, $29$, $52$ and $63$ per cent of
sessions at $0.05$, $0.1$, $0.2$ and $0.5$\,mag.

\emph{Phase curves.} A surface that shines with its own light, or that reflects as a
smooth specular body does, brightens little with decreasing phase angle. The ATLAS
phase curves in two bands~\cite{Mahlke2021} give $39239$ fits of the $H$, $G_1$, $G_2$
function~\cite{Muinonen2010} spanning phase angles below $7^\circ$ and above $25^\circ$,
with a median slope of $0.0277$\,mag per degree from $5$ to $30^\circ$. The $221$ curves
pinned at the flattest slope the function allows, $0.0179$\,mag per degree, have median
residuals of $0.28$\,mag, rotation swamping the phase signal; the $107$ of them with
measured albedos have $0.03$ to $0.74$, none above the $0.8$ that a flat curve on a bright
surface would need to stand out.

\emph{Beaming.} A body heated by more than the Sun shows a hotter dayside than a rough
sunlit surface can, a NEATM~\cite{Harris1998} beaming parameter below about $0.6$. Of
the $74609$ NEOWISE thermal fits with a free beaming parameter~\cite{Mainzer2011}, of
$68500$ objects, eight fall below $0.6$ at $2\sigma$, each at a single epoch, and the
lowest sits at the fit's floor of $0.3$; the population's first percentile is $0.68$,
and no object stays below $0.6$ at every epoch it was fitted.

% =============================================================================
\section{Shapes and spins}\label{sec:shapes}
% =============================================================================

\emph{Plates.} A collisional fragment's lightcurve amplitude rarely exceeds $1.2$\,mag at
low phase, while a plate of axis ratio $10$ reaches $2.5$\,mag. In the LCDB's detail
records, $31008$ of quality $2$ or $3$ at phase angles up to $15^\circ$, with amplitudes
reduced to zero phase at $0.03$\,mag per degree, five records exceed $1.6$\,mag, $0.02$
per cent against the $0.11$ per cent of the quality-1 records, one exceeds $2.0$ and none
$2.5$. All five are single TESS records~\cite{Pal2020} of slow rotators of $9$ to $627$\,h,
three marked ambiguous; no body exceeds a $4:1$ axis ratio at low phase.

\emph{Straight limbs.} A stellar occultation traces a body's limb chord by chord. The
occultation bundle of the Small Bodies Node holds $395$ events of fit quality $3$ or $4$
with five or more positive chords from distinct sites, of which $376$ are measured once
$19$ whose chords disagree or whose fitted ellipse is far from the body's size are set
aside. Each chord's ends are placed on the fundamental plane, and the longest run of
four or more consecutive limb points that a
line fits within the timing errors is measured against the diameter and against the
fitted ellipse's points at the same angles and errors. The median run is $0.21$ of the
diameter and the $90$th percentile $0.67$. Twenty-five runs exceed the ellipse's null at
$p<0.01$, the longest on (3317)~Paris in 2020 ($0.97$ of the diameter over six sites),
(134)~Sophrosyne in 1980 ($0.86$), (64)~Angelina in 2023 ($0.82$), (39)~Laetitia in 1998
($0.79$) and (2)~Pallas in 1983 ($0.70$ over six of $120$ chords). Their shape
models~\cite{Durech2010} carry the flat facets that an ellipse lacks, so the null that
decides is the models projected at each event's geometry.

\emph{Spin strength.} A gravitational aggregate of $2000$\,kg\,m$^{-3}$ flies apart
faster than about $2.3$\,h~\cite{Pravec2000}, and a faster spin needs cohesion that grows
with size~\cite{Holsapple2007}. Of the $8854$ asteroids with periods and diameters in the
LCDB, no kilometre body spins faster than ten minutes or needs $100$\,kPa; the most
demanding are 2019~MA2, $0.49$\,km across from its $H$ and spinning in $2.62$\,min at
quality~$2$, which needs $97$\,kPa, within a monolith's strength, and 2008~RA106,
$3.7$\,km and $54$\,min at quality $2-$, which needs $11$\,kPa. The secure fast rotators
of the $31149$ periods of quality $2$ or $3$ are the cohesive ones already published.

\emph{Replicated bodies.} Objects made to one design would share a period, a size and a
lightcurve amplitude on orbits that collisions do not relate. Among $4315$ numbered
asteroids with LCDB periods of quality~$3$ and the proper elements of the asteroid-family
bundle~\cite{Nesvorny2015}, pairs that agree in period to $2\times10^{-4}$, in $H$ to
$0.1$\,mag and in amplitude to $0.05$\,mag while more than $300$\,m\,s$^{-1}$ apart in
the Zappal\`a metric~\cite{Zappala1990} number $14$, against $8.6\pm2.8$ with the periods
shuffled within half-magnitude bins, and one group of three forms, against
$0.04\pm0.18$, once in $25$ shuffles: (3305)~Ceadams, (5122)~Mucha and (2019)~van Albada,
at $2.7290$ to $2.7294$\,h, $H=12.3$ to $12.4$ and amplitudes $0.20$ to $0.23$, in three
families. Their Gaia reflectance spectra are all S-type, with red slopes to $1.21$ to
$1.26$ at $770$\,nm and the $1\,\mu$m band, and their period lies in the pile-up of fast
rotators near the spin barrier, where periods crowd.

% =============================================================================
\section{Orbits}\label{sec:orbits}
% =============================================================================

\emph{Area-to-mass ratios.} A radial non-gravitational term $A_1$ measures the push of
sunlight or outgassing per unit mass, and for a body pushed by sunlight alone gives its
area-to-mass ratio. Of the $301$ objects with $A_1$ in JPL's database, three are lighter
than a rubble pile under the most favourable assumptions: the active asteroid
362P/(2008~GO98), the interstellar object 1I/2017~U1 and the dark comet
2016~NJ33~\cite{Seligman2023}, two of them with transverse or normal terms that mark
activity.

\emph{Unreachable orbits.} The delivery model of Granvik et al.~\cite{Granvik2018} gives
the density of near-Earth objects from each natural source over orbital elements, and
leaves cells empty that no source reaches. The JPL database holds $37$ near-Earth objects
in those cells: $24$ retrograde, all at eccentricities of $0.74$ or more, the
Halley-type comets and comet-like asteroids; none retrograde below $0.3$; one
near-circular orbit inside $0.75$\,au above $25^\circ$, 2025~LP, within the Atira
population; and $12$ above $60^\circ$ with eccentricities below $0.25$, which Kozai
cycles~\cite{Kozai1962} reach from lower inclinations.

\emph{Hyperbolic meteor streams.} Meteoroids from another star, or debris shed by
something travelling between stars, would arrive on one incoming asymptote at one
speed. The Global Meteor Network~\cite{Vida2021} holds $3.58$ million orbits from 2018
to 2026, $40753$ sporadic meteors hyperbolic at $3\sigma$; each orbit gives its incoming
asymptote and excess speed, and a meteor's neighbours within $1^\circ$ and
$3$\,\kms\ are counted against the density within $10^\circ$. Thirty-five meteors fall
below $p=10^{-4}$ in ten groups. Eight belong to IAU showers~\cite{Jopek2017} whose
speeds are overestimated (Perseids, Orionids, $\eta$~Aquariids and the $\xi$~Perseids).
The largest, $16$ meteors from 2022 to 2026 at solar longitudes $80$ to $214^\circ$ on
one asymptote at ecliptic longitude $134.5^\circ$ and latitude $-1.6^\circ$ at
$12$\,\kms, is matched by a $16$-member group in one of three realisations of a $1.5$ per
cent speed overestimate shared by the elliptic sporadics, $74000$ simulated meteors, and
fifteen realisations thinned to the observed number at $1.0$ to $2.0$ per cent give $3$
to $31$ meteors below $p=10^{-4}$ and groups of up to $14$. SonotaCo's independent
$353231$ orbits~\cite{SonotaCo2009} hold one meteor on that asymptote against $0.16$ from
the local density and about $1.6$ that a real stream would give. A stream would also
spread along the loci on which Earth's gravity focuses it, reaching $\pm52^\circ$ of
latitude, half of them south, as Earth crosses the asymptote near solar longitude
$134.5^\circ$: the August 3 to 11 loci hold $529$ meteors against a mean of $16$ on $11$
control asymptotes, but $510$ of them lie north, where the Perseids ($204$ tagged) and
the apex sporadics are, and $19$ south, against $10$ by chance. The group is the speed
bias acting on the shower population, and no star within $100$\,pc shares its velocity
on its line.

% =============================================================================
\section{Lunar flashes}\label{sec:moon}
% =============================================================================

Impacts on the Moon's night side flash thermally for tens of milliseconds; a flash from
a surface source, a lamp or a laser, would be hotter or cooler than the impact plume,
longer, or recur at one place. NELIOTA's flash lists~\cite{Liakos2020,Liakos2024} hold
$284$ flashes from 2017 to 2023, $188$ of them validated and $192$ with $R$ and $I$
photometry, located to $0.5^\circ$. Colour temperatures from $R-I$ through the Planck
ratio, with Vega's colour as the zero point, have a median of $2935$\,K and a $5$th to
$95$th percentile range of $2027$ to $5020$\,K; none is hotter than $8000$ or cooler than
$1000$\,K at $3\sigma$, and none lasts longer than $0.5$\,s, the longest $429$\,ms. Pairs
within $1.5^\circ$ number $21$ against $23.9\pm4.9$ with the longitudes permuted among the
flashes ($p=0.75$); the closest, two validated flashes $0.10^\circ$ apart on consecutive
Perseid nights of 2018 August at $(+26.7^\circ,+60.2^\circ)$, with temperatures of $4100$
and $3000$\,K, is one that a tenth of a pair is expected to match.

\begin{table}[t]
\centering\small
\caption{The thirteen tests, their samples and results.}
\label{tab:tests}
\begin{tabular}{>{\raggedright\arraybackslash}p{0.2\textwidth} >{\raggedright\arraybackslash}p{0.31\textwidth} >{\raggedright\arraybackslash}p{0.41\textwidth}}
\toprule
Test & Sample & Result \\
\midrule
Radar albedos & $315$ radar detections & albedos above $0.3$ are M, X, E and one Sq type; no specular reflector \\
Reflectance spectra & $58208$ Gaia DR3 spectra & no photovoltaic or metal surface; $88$ per cent recovery at $10$ per cent of the light \\
Glints & $4723$ ALCDEF sessions & no spikes recurring at the rotation period at $p<0.01$ \\
Phase curves & $39239$ ATLAS curves & the flattest are rotation-dominated, albedos $0.03$ to $0.74$ \\
Beaming & $74609$ NEOWISE fits & eight single-epoch fits below $\eta=0.6$; none at every epoch \\
Plates & $31008$ LCDB records & none above $2.5$\,mag at zero phase \\
Straight limbs & $395$ occultations & longest runs on bodies with faceted shape models \\
Spin strength & $8854$ asteroids & no kilometre body under $10$\,min or above $100$\,kPa \\
Replicated bodies & $4315$ asteroids & one S-type triple at the spin barrier, one in $25$ \\
Area-to-mass & $301$ radial terms & three light bodies, each explained \\
Unreachable orbits & $37$ near-Earth objects & comets, an Atira and Kozai-reachable orbits \\
Hyperbolic streams & $40753$ hyperbolic meteors & groups matched by a $1.5$ per cent speed bias; no focal line \\
Lunar flashes & $284$ NELIOTA flashes & thermal, $2000$ to $5000$\,K; recurrence at chance \\
\bottomrule
\end{tabular}
\end{table}

% =============================================================================
\section{Related work}\label{sec:related}
% =============================================================================

Freitas~\cite{Freitas1983} set out the search for extraterrestrial artifacts in the
Solar System, Papagiannis~\cite{Papagiannis1983} argued for exploring the asteroid belt
for them, Haqq-Misra and Kopparapu~\cite{HaqqMisra2012} estimated how little of the
Solar System has been searched at the sizes a probe would have, and
Benford~\cite{Benford2019} proposed the co-orbital objects of the Earth as the likeliest
places for one. The tests here rest on
the archives' own studies: the radar survey of metal-rich asteroids~\cite{Shepard2015},
the Gaia reflectance spectra~\cite{Galluccio2023}, the ATLAS phase
curves~\cite{Mahlke2021}, the spin barrier and its cohesive
exceptions~\cite{Pravec2000,Holsapple2007}, the dark comets~\cite{Seligman2023}, the
debiased near-Earth population~\cite{Granvik2018} and the meteor networks' orbit
catalogues~\cite{Vida2021,SonotaCo2009}, which this paper puts to the use those
proposals describe.

% =============================================================================
\section{Conclusion}\label{sec:conclusion}
% =============================================================================

No test finds an artifact among the small bodies. Every bright radar echo is a metallic
or enstatite asteroid, no reflectance spectrum carries a photovoltaic surface, no
lightcurve glints at its rotation period, no phase curve or thermal fit needs a source
of light other than the Sun, no amplitude, limb or spin exceeds what natural shapes and
strengths allow, the one coincidence of period, size and amplitude is the expected
crowding at the spin barrier, every light body is explained, no near-Earth orbit lies
beyond the natural sources' reach, the hyperbolic meteor groups are speed biases, and the
lunar flashes are impacts.

The reach differs from test to test. The spectra exclude photovoltaic cover of a tenth of
a body's light among bodies Gaia measured, down to about $G=20$, and say nothing of
silver or aluminium; the glint search reaches spikes of $0.1$\,mag in a third of the
sessions; the radar sample is the $315$ bodies radar has reached; the occultation test
needs the shape-model null to separate a made facet from a natural one; and an object
smaller or fainter than these surveys' limits, or one that keeps its distance from the
inner Solar System, is not constrained.

\section*{Acknowledgements}

This work uses the Planetary Data System Small Bodies Node (the asteroid radar
compilation, the ALCDEF photometry, the Asteroid Lightcurve Database, the occultation
bundle and the asteroid-family bundle), Gaia DR3 through CDS and ESA, the ATLAS phase
curves through VizieR, the NEOWISE diameters and albedos from the NASA/IPAC Infrared
Science Archive, JPL's Small-Body Database, the Global Meteor Network's trajectory
summaries, SonotaCo's orbit catalogue, the IAU Meteor Data Center, the NELIOTA flash
tables, the optical constants of refractiveindex.info, and the software NumPy, SciPy and
astropy.

\begin{thebibliography}{99}\small\raggedright
\setlength{\itemsep}{1pt}

\bibitem{Freitas1983}
R.~A. Freitas, The search for extraterrestrial artifacts (SETA), \emph{J. Br. Interplanet.
Soc.} \textbf{36}, 501 (1983).

\bibitem{Papagiannis1983}
M.~D. Papagiannis, The importance of exploring the asteroid belt, \emph{Acta Astronaut.}
\textbf{10}, 709 (1983).

\bibitem{HaqqMisra2012}
J.~Haqq-Misra and R.~K. Kopparapu, On the likelihood of non-terrestrial artifacts in the
Solar System, \emph{Acta Astronaut.} \textbf{72}, 15 (2012).

\bibitem{Benford2019}
J.~Benford, Looking for lurkers: co-orbiters as SETI observables, \emph{Astron. J.}
\textbf{158}, 150 (2019).

\bibitem{Shepard2015}
M.~K. Shepard, P.~A. Taylor, M.~C. Nolan, et al., A radar survey of M- and X-class
asteroids. III. Insights into their composition, hydration state, and structure,
\emph{Icarus} \textbf{245}, 38 (2015).

\bibitem{Ostro1991}
S.~J. Ostro, D.~B. Campbell, J.~F. Chandler, et al., Asteroid 1986 DA: radar evidence for
a metallic composition, \emph{Science} \textbf{252}, 1399 (1991).

\bibitem{Galluccio2023}
Gaia Collaboration, L.~Galluccio, M.~Delbo, et al., Gaia Data Release 3: reflectance
spectra of Solar System small bodies, \emph{Astron. Astrophys.} \textbf{674}, A35 (2023).

\bibitem{Polyanskiy2024}
M.~N. Polyanskiy, Refractiveindex.info database of optical constants, \emph{Sci. Data}
\textbf{11}, 94 (2024).

\bibitem{Warner2009}
B.~D. Warner, A.~W. Harris and P.~Pravec, The asteroid lightcurve database,
\emph{Icarus} \textbf{202}, 134 (2009).

\bibitem{Mahlke2021}
M.~Mahlke, B.~Carry and L.~Denneau, Asteroid phase curves from ATLAS dual-band
photometry, \emph{Icarus} \textbf{354}, 114094 (2021).

\bibitem{Muinonen2010}
K.~Muinonen, I.~N. Belskaya, A.~Cellino, et al., A three-parameter magnitude phase
function for asteroids, \emph{Icarus} \textbf{209}, 542 (2010).

\bibitem{Harris1998}
A.~W. Harris, A thermal model for near-Earth asteroids, \emph{Icarus} \textbf{131}, 291
(1998).

\bibitem{Mainzer2011}
A.~Mainzer, T.~Grav, J.~Masiero, et al., Thermal model calibration for minor planets
observed with Wide-field Infrared Survey Explorer/NEOWISE, \emph{Astrophys. J.}
\textbf{736}, 100 (2011).

\bibitem{Pal2020}
A.~P\'al, R.~Szak\'ats, Cs.~Kiss, et al., Solar System objects observed with TESS: first
data release, bright main-belt and Trojan asteroids from the southern survey,
\emph{Astrophys. J. Suppl.} \textbf{247}, 26 (2020).

\bibitem{Durech2010}
J.~\v{D}urech, V.~Sidorin and M.~Kaasalainen, DAMIT: a database of asteroid models,
\emph{Astron. Astrophys.} \textbf{513}, A46 (2010).

\bibitem{Pravec2000}
P.~Pravec and A.~W. Harris, Fast and slow rotation of asteroids, \emph{Icarus}
\textbf{148}, 12 (2000).

\bibitem{Holsapple2007}
K.~A. Holsapple, Spin limits of Solar System bodies: from the small fast-rotators to
2003 EL61, \emph{Icarus} \textbf{187}, 500 (2007).

\bibitem{Nesvorny2015}
D.~Nesvorn\'y, M.~Bro\v{z} and V.~Carruba, Identification and dynamical properties of
asteroid families, in \emph{Asteroids IV}, ed. P.~Michel, F.~E. DeMeo and W.~F. Bottke
(Univ. Arizona Press, 2015), p.~297.

\bibitem{Zappala1990}
V.~Zappal\`a, A.~Cellino, P.~Farinella and Z.~Kne\v{z}evi\'c, Asteroid families. I.
Identification by hierarchical clustering and reliability assessment, \emph{Astron. J.}
\textbf{100}, 2030 (1990).

\bibitem{Seligman2023}
D.~Z. Seligman, D.~Farnocchia, M.~Micheli, et al., Dark comets? Unexpectedly large
nongravitational accelerations on a sample of small asteroids, \emph{Planet. Sci. J.}
\textbf{4}, 35 (2023).

\bibitem{Granvik2018}
M.~Granvik, A.~Morbidelli, R.~Jedicke, et al., Debiased orbit and absolute-magnitude
distributions for near-Earth objects, \emph{Icarus} \textbf{312}, 181 (2018).

\bibitem{Kozai1962}
Y.~Kozai, Secular perturbations of asteroids with high inclination and eccentricity,
\emph{Astron. J.} \textbf{67}, 591 (1962).

\bibitem{Vida2021}
D.~Vida, D.~\v{S}egon, P.~S. Gural, et al., The Global Meteor Network: methodology and
first results, \emph{Mon. Not. R. Astron. Soc.} \textbf{506}, 5046 (2021).

\bibitem{Jopek2017}
T.~J. Jopek and Z.~Ka\v{n}uchov\'a, IAU Meteor Data Center: the shower database, a status
report, \emph{Planet. Space Sci.} \textbf{143}, 3 (2017).

\bibitem{SonotaCo2009}
SonotaCo, A meteor shower catalog based on video observations in 2007--2008,
\emph{WGN} \textbf{37}, 55 (2009).

\bibitem{Liakos2020}
A.~Liakos, A.~Z. Bonanos, E.~M. Xilouris, et al., NELIOTA: methods, statistics, and
results for meteoroids impacting the Moon, \emph{Astron. Astrophys.} \textbf{633}, A112
(2020).

\bibitem{Liakos2024}
A.~Liakos, A.~Z. Bonanos, E.~M. Xilouris, et al., NELIOTA: new results and updated
statistics after 6.5 years of lunar impact flashes monitoring, \emph{Astron. Astrophys.}
\textbf{687}, A14 (2024).

\end{thebibliography}

\end{document}
